% CONSUMO ENERGESTICO Y ESTADO ACTUAL

Bajo el paradigma actual de la computación distribuida pueden distinguirse dos modelos principales de despliegue en la nube \cite{cloud-bases}: las nubes públicas y las nubes privadas. Las primeras ofrecen recursos de computación como un servicio accesible para terceros, permitiendo a los clientes utilizar la capacidad de procesamiento de grandes Centros de Procesamiento de Datos (CPDs) sin necesidad de gestionar ni mantener la infraestructura subyacente. Por el contrario, las nubes privadas son administradas por la propia organización sobre una infraestructura dedicada, lo que implica asumir costes adicionales asociados al mantenimiento del hardware, la administración de los sistemas y el consumo energético necesario para garantizar su funcionamiento.\

La creciente adopción de estos modelos de computación ha impulsado la construcción de grandes infraestructuras capaces de atender millones de peticiones de forma simultánea. Entre ellas destacan las Redes de Distribución de Contenido (CDN, Content Delivery Networks), cuya finalidad es distribuir contenido multimedia a través de redes de alta velocidad para reducir la latencia y mejorar la experiencia de los usuarios. Tanto los CPDs como las CDN requieren una gran cantidad de recursos energéticos para mantener un funcionamiento continuo las 24 horas del día, los 365 días del año \cite{cpds-consumption-energy}. Este consumo no se limita únicamente a los equipos de procesamiento, sino que también incluye los sistemas de refrigeración necesarios para mantener unas condiciones de operación adecuadas, siendo el agua uno de los recursos más utilizados en la industria para este fin\

El incremento del número de CPDs y de la capacidad de procesamiento requerida ha supuesto un importante impacto desde el punto de vista medioambiental. Según el informe publicado por la International Energy Agency (IEA) \cite{iae-report-consumption}, la demanda energética asociada a los centros de datos continúa aumentando debido al crecimiento de las cargas de trabajo digitales, lo que ha impulsado la necesidad de ampliar la capacidad de generación eléctrica \cite{iae-report-generation}. Entre los principales factores responsables de este incremento destacan el auge de la inteligencia artificial, el procesamiento masivo de datos y la construcción de nuevas infraestructuras tecnológicas, como centros de datos y redes de distribución de contenido, que requieren un suministro energético continuo y sistemas de refrigeración de gran capacidad.\

Como consecuencia de esta elevada demanda, numerosos proyectos tecnológicos se enfrentan actualmente a limitaciones en el acceso a la red eléctrica. El informe de la IEA indica que existen solicitudes de conexión que representan aproximadamente 2.500 GW de nueva demanda eléctrica, muchas de las cuales permanecen en espera debido a la insuficiente capacidad de generación y distribución energética. Aunque el propio informe señala que parte de estas solicitudes corresponden a proyectos que finalmente podrían no materializarse, pone de manifiesto que el crecimiento de la demanda energética constituye un desafío real y de alcance global.\\

\subsection{Generacion Energetica}
Como respuesta al incremento de la demanda energética y a la necesidad de reducir el impacto medioambiental, durante los últimos años se han producido cambios significativos tanto en las fuentes de generación eléctrica como en su distribución. Según el informe publicado por la International Energy Agency (IEA) \cite{iae-report-generation}, la tendencia a medio y largo plazo apunta hacia un aumento progresivo de la participación de fuentes de energía renovables y de tecnologías de bajas emisiones. Entre ellas, la energía solar fotovoltaica destaca como la tecnología con mayor crecimiento a nivel mundial, impulsada principalmente por las inversiones realizadas por China y Estados Unidos durante los últimos años. Este crecimiento ha permitido reducir la dependencia del carbón en numerosos sistemas eléctricos, incluso en un contexto de aumento de la demanda energética. Asimismo, la biomasa continúa consolidándose como una alternativa renovable relevante debido al aprovechamiento energético de residuos, mientras que la energía nuclear ha alcanzado niveles históricos de generación eléctrica, contribuyendo a satisfacer el incremento de la demanda mundial.\

No obstante, los combustibles fósiles continúan desempeñando un papel predominante en el sistema energético global. Aproximadamente el 35\% de la energía consumida procede del petróleo y alrededor del 20\% del gas natural, que continúa siendo una de las principales fuentes de generación eléctrica a nivel mundial. Aunque el incremento de las energías renovables y las mejoras en eficiencia energética han permitido limitar el crecimiento de las emisiones de dióxido de carbono, las emisiones globales alcanzaron aproximadamente 38,400 millones de toneladas de CO$_2$, impulsadas principalmente por el uso de carbón, gas natural y diversos procesos industriales \cite{iae-report-generation}.

Ante este escenario, la expansión de las fuentes de generación con bajas emisiones constituye uno de los principales objetivos de las políticas energéticas internacionales. En este sentido, la COP \cite{COP2030} establece como uno de sus objetivos incrementar de forma significativa la participación de las energías renovables y de la energía nuclear en el mix energético antes de 2030, al mismo tiempo que se continúa reduciendo la dependencia del carbón. Para alcanzar estas metas también será necesario modernizar las infraestructuras eléctricas mediante el aumento de la capacidad de transporte, almacenamiento y distribución de energía. Estas medidas permitirán integrar una mayor proporción de energías renovables, satisfacer el crecimiento de la demanda eléctrica y contribuir a una reducción progresiva de las emisiones globales de gases de efecto invernadero.


\subsection{Eficiencia Energetica}

% - Esta reudccion del consumo energetico se deberse a varios factores, siendo algunos de ellos debido al "green users" y el conjunto de mejoras en multiples procesos,
%  que han hecho mas eficinete diferentes partes de la cadena de  consumo de contenido. Desde el dispositovo final, pasando por los protocolos de delivery de contenido hasta los metodos de codificacion. 